\section{Representative activation and the limits of forest selection}
\label{sec:representative-activation}

Report~115 replaces exceptional-row order enumeration with a
single-exponential representative-family search for a fixed local
cancellation problem.  This 12 October continuation reviews that argument
and executes its certificates against the maintained native geometry.
The local theorem is useful, but our native experiments find no additional
core coverage over the existing maximum-weight top forest.  More
significantly, they give exact obstructions to resolving two outstanding
cores by changing either forest.  We retain the selector as research
tooling rather than add an unmeasured stage to the default recognizer.

\subsection{The exact local problem}

Fix the source handle structure and initially a primal zero/one forest.
Let $E$ index its two-handles, and let $N=(n_{eh})$ be the exact
\emph{unsigned} intersection counts on one-handles outside that forest.
The dual graph $D$ has one boundary root and the interior three-handles
as vertices; its edge $e$ retains the two endpoint occurrences of
two-handle $e$.  Loops and parallel edges are permitted.  Write $m=|E|$
and $r=|V(D)|-1$.  The source dual graph is connected to its boundary root.

A spanning tree $T$ deletes the corresponding two-handles by root-outward
two/three cancellations.  The remaining rows must then peel: row $e$ can
be deleted on column $h$ exactly when $n_{eh}=1$ and every other live row
has zero in that column.  Used columns may be retired.  Keeping their
implications after use is harmless, since all their support rows have
already disappeared.  Additional initial row deletions can only help.
Consequently, an independent set of dual edges whose deletion permits
complete peeling extends to a successful spanning tree.

Counts are never reduced modulo a field in this criterion.  For example,
the matrix with rows $(2,1)$ and $(1,2)$ is a stopping matrix, whereas
its reduction modulo two peels.  Finite-field calculations below apply
only to the incidence representation of the \emph{dual graph}.

\subsection{Reducing peeling to a small activation boundary}

A column is usable if it has a unit entry.  Unit singletons are anchors.
A usable support-two column with rows $u,v$ gives the implication
$u\longrightarrow v$ when $n_{vh}=1$, and conversely when $n_{uh}=1$.
Let $G$ be this directed graph on rows, and write $\operatorname{Reach}$
for its reflexive reachability relation.  Define $B$ as the union of
\emph{all} support rows of usable columns having support at least three,
and put $b=|B|$.  Heavy entries remain prerequisites and therefore remain
in $B$.  Columns with no unit entry never supply a pivot.

A source strongly connected component of $G$ containing neither an
anchor nor a row in $B$ is mandatory.  Denote these disjoint components
by $U_1,\ldots,U_q$.  Every successful initial deletion set must meet
each $U_i$: its first deleted row cannot arise from an anchor, an incoming
low-column implication, or a high-column pivot.  Conversely, any one
seed in $U_i$ activates that entire component and its descendants.

Set $H(e)=B\cap\operatorname{Reach}(e)$.  On subsets of $B$, let
$\operatorname{cl}$ saturate the following rules: include the reachability
from anchors and all mandatory components; propagate from any activated
$x\in B$ to $H(x)$; and activate a unit target of a high column whenever
every other support row of that column is activated.  Mandatory components
are included symbolically only on the condition that actual independent
representatives will be supplied.

\begin{proposition}[Exact activation state]
After choosing one representative from every mandatory component and an
additional seed set $X$, the deleted exceptional rows are exactly
$\operatorname{cl}(\bigcup_{e\in X}H(e))$.  All rows peel if and only
if this set equals $B$.
\end{proposition}
\begin{proof}
Each closure rule is a valid low-column path or a valid private high-column
pivot.  Conversely, inspect a literal peeling trace: low moves give
reachability, and a high move has its target and all prerequisites in $B$.
Thus every move is represented by the closure.  Once all of $B$ disappears,
each source component of $G$ contains an activated anchor, a mandatory seed,
or an activated exceptional row.  Every component lies downstream of a
source component, so low propagation deletes the remaining rows.
\end{proof}

Every successful tree has a witness using at most $b$ additional seeds.
Choose its mandatory representatives first, then repeatedly take a tree
edge whose influence increases the closed state.  If this stops before
$B$, the entire tree cannot activate $B$, a contradiction.  Each useful
edge adds a bit, so at most $b$ are needed.  Their union with the mandatory
representatives remains a forest.  This replaces an order of all
exceptional rows by an activation state and a small independent seed set.

\subsection{Contraction, representatives, and the probability guarantee}

Represent $D$ by reduced oriented incidence columns $v_e\in\mathbb F_p^r$.
Independent columns are precisely forests, over every prime field; forest
minors have determinant $\pm1$.  For each mandatory component introduce
independent variables and the formal column
$y_i=\sum_{e\in U_i}z_{i,e}v_e$.  A minor of $[Y,V_X]$ is a polynomial
whose distinct representative choices have distinct monomials.
It is nonzero exactly when $X$ extends by one representative from each
$U_i$ to a forest.  A nonzero specialized minor therefore guarantees a
genuine extension.  Specialization can lose witnesses but cannot invent
one.

Sample the coefficients independently in $\mathbb F_p$.  Reject a round
if $Y$ loses rank $q$; otherwise quotient by its span and apply an
independent random projection to dimension $d=\min(b,r-q)$.  Search these
projected vectors.  For a fixed feasible $k$-seed witness, failure of a
nonzero augmented minor has probability at most $q/p$.  Conditional on
surviving, failure of a projected $k$-minor has probability at most $k/p$.
Thus a fixed witness survives with probability at least $1-(q+k)/p$.
Choosing $p>2(q+b)$ gives success probability greater than one half per
complete independent round.  There is no union bound over all witnesses.

At state $(S,k)$ retain actual $k$-seed sets whose exterior vectors form
a basis of the span of all candidates at that state.  There are at most
$\binom dk$ such representatives.  Extend the \emph{original decomposable}
vectors, not their Gaussian-eliminated linear combinations.  For a future
set $Z$, the map $w\mapsto w\wedge w_Z$ is linear.  If it is nonzero on
any discarded vector, it is nonzero on some retained vector.  The future
activation depends only on $S$ and $Z$, so that representative preserves
the same successful extension.  This is the classical linear-matroid
representative-family mechanism \cite{activation-fomin}; the application
here rests on the exact activation-state proposition above.

There are at most $2^b$ masks and
$m2^b\sum_k\binom dk=m2^{b+d}$ candidate extensions.  Straightforward
exterior multiplication and Gaussian insertion give the conservative
field-operation bound
\[
  2^b8^d\operatorname{poly}(\Lambda)\leq
  16^b\operatorname{poly}(\Lambda),
\]
where $\Lambda$ is the explicit encoded matrix/graph size.  Storage is
bounded by $2^b4^d\operatorname{poly}(\Lambda)$.  The field has
$O(\log(q+b+2))$-bit elements; all graph, closure, quotient and certificate
work contributes polynomial factors.  With $\lambda$ complete independent
rounds, positive miss probability is at most $2^{-\lambda}$.  Caps or
fixed reproducibility seeds do not inherit that guarantee.

When $q=0$ and $r\leq b$, the original binary graph representation already
has suitable dimension and is faithful.  The search is then deterministic.
A negative frontier packet lists basis witnesses at each nonterminal state.
The independent checker reconstructs exterior coordinates by binary graph
minors and activation by literal peeling.  It checks the initial empty
witness and span coverage of every independent strict-progress extension.
Absence of the terminal state then proves local infeasibility.  In the
general sampled regime an exhausted search says UNKNOWN.

\subsection{Native authentication and the all-forest obstruction}

The adapter emits the established
\path{core-three-handle-reduction-v1} schema.  Our audit first authenticates
the full source parallelity packet, counts raw cycle flags without
saturation, and preserves both dual endpoint occurrences.  The unchanged
native replay derives owners from source vertex-link triangles and normal
coordinates.  After it accepts a proposal and renumbers the surviving
rows, we build and independently check a \emph{new} core certificate.
No old basis-dependent port words are reused.  These experiments stop at
the core; they do not assert additional whole-cut certificates.

There is a useful stronger negative test.  Instead of deleting the columns
of a chosen lower forest, retain all source one-handle columns and construct
the mandatory components on this full matrix.  A successful fixed-lower
trace would also be a trace on the full matrix: its chosen pivot columns
and row counts are identical, and extra unused columns cannot obstruct
a private-column move.  Hence failure for every dual tree on the full
matrix rules out every pair of lower and upper forests for this method.

\begin{proposition}[Mandatory-rank obstruction for all forest pairs]
If a collection $J$ of mandatory source components of the full matrix
satisfies
\[
 |J|>\operatorname{rank}_D\left(\bigcup_{j\in J}U_j\right),
\]
no choice of a dual spanning tree and a primal spanning forest gives
complete private-column cancellation.
\end{proposition}
\begin{proof}
Every successful deletion tree must meet each of the disjoint mandatory
components.  Selecting one of its edges per component gives $|J|$
independent edges in their union, contrary to the displayed rank inequality.
Thus no top tree succeeds on the full matrix.  The preceding monotonicity
argument transfers this obstruction to every contracted lower forest.
\end{proof}

This does not obstruct handle slides, changes to the source decomposition,
the Euler/free-group terminal, or handlebody structure itself.  In particular
it is never a knottedness certificate.  In the two measured examples even
the elementary bound $q>r$ suffices:
\begin{center}
\begin{tabular}{lrrrr}
source & rows & columns & mandatory $q$ & dual rank $r$\\
\hline
\path{nested_1_4_6/interior_links} &48&98&7&6\\
\path{interior_loop/interior_links} &10&21&2&1
\end{tabular}
\end{center}
Both have $b=0$.  Their exact component lists, matrix digests and independently
replayed rank certificates are retained with the native audit.  This sharpens
the earlier fixed-top residual observations: neither case can be repaired
by further forest selection within this cancellation model.

\subsection{Executed evidence and integration decision}

On a scratch copy, all 71 delivered tests pass.  The full supplied algebra
audit reproduces 8,061 exhaustive/random matrix-graph inputs: 3,117 positive
answers, 3,391 exact local negatives and 1,553 UNKNOWN negatives, with no
missed positive in that finite sample.  Another 300 planted positive inputs
succeed in 318 rounds.  Additional comparisons cover 16,038 forced matroid
queries, 2,369 closure queries, 1,951 sampled-independence checks, 255
projected-search instances and all 503 assignments in four small-field
experiments.  These remain algebraic validation, distinct from source
realizability or a probability proof.

The native audit covers 126 fully authenticated cuts: the existing 120 plus
the six signed-annulus fixtures.  Only 28 have interior three-handles.
For each cut it tests a fixed source lower forest, a fixed maximum-passage
forest, and the full-column/deferred variant.  Every positive proposal and
rebuilt core certificate independently replays.  The comparison is:
\begin{center}
\begin{tabular}{lrrr}
lower policy & current positives & selector positives & exact local refusals\\
\hline
source &118&118&8\\
maximum passage &124&124&2\\
deferred &124&124&2
\end{tabular}
\end{center}
Of the 378 trials, 328 succeed on the selector's ordinary incumbent tree,
36 use the empty-input route, two use faithful activation, six produce
faithful-frontier obstructions and six produce mandatory-rank obstructions.
No random projection or cap exhaustion occurs.  The boundary parameter
reaches 47, but those large-boundary trials succeed before the cap of 20.
The transition allowance is 100,000 and the declared sampling allowance
is eight rounds; these finite settings do not amount to an uncapped theorem
test.  The two faithful-activation cases are the loop examples already
resolved by the maintained maximum-weight tree.

A separate deterministic search applies 30 relabellings to each of six
sources, giving 180 relabelled cuts and 540 lower-policy trials.  All
maximum-passage and deferred trials already succeed.  Of the fixed-source
trials, 178 succeed and two receive exact local obstructions.  There is no
new positive.  This search uses geometric preparation; its planned full
source authentication for new hits is never triggered.  It is therefore
not counted as another 180 fully authenticated source packets.

No timings here establish a native speedup, and no maintained runtime path
changes.  The original report is preserved.  Reproducible drivers are in
\path{fast/activation_forest_research}; outputs, source hashes and local
certificates are in \path{synthesis/data/activation-forest-*}.
For a global quasi-polynomial consequence one would need bounds on the
exceptional parameter, source sizes, useful-surface discovery and the
number and repair cost of hierarchy calls, together with the required
marked successor geometry.  The local $2^{O(b)}$ theorem and the all-forest
obstructions supply neither those bounds nor that geometry.
